% ============================================================================
% METHODS: POLYPHARMACY RISK ASSESSMENT AND ALTERNATIVE THERAPY SUGGESTION
% Publication Format - Methods Section
% ============================================================================

\documentclass[11pt,a4paper]{article}
\usepackage[utf8]{inputenc}
\usepackage{amsmath,amssymb}
\usepackage{booktabs}
\usepackage{geometry}
\usepackage{graphicx}
\usepackage{hyperref}
\usepackage{float}
\usepackage{xcolor}
\geometry{margin=1in}

\hypersetup{
    colorlinks=true,
    linkcolor=blue,
    filecolor=magenta,
    urlcolor=cyan,
}

\title{\textbf{Methods}\\
\large Polypharmacy Risk Assessment and Alternative Therapy Suggestion}
\author{Drug-Drug Interaction Knowledge Graph Project}
\date{}

\begin{document}

\maketitle

\tableofcontents
\newpage

% ============================================================================
\section{Polypharmacy Risk Assessment}
% ============================================================================

\subsection{Polypharmacy Risk Index (PRI)}

We developed the Polypharmacy Risk Index (PRI) to quantify individual drug risk within multi-drug regimens. The PRI integrates four network-based metrics weighted according to their clinical relevance:

\begin{equation}
\text{PRI}(d) = 0.25 \cdot C_{\text{degree}}(d) + 0.30 \cdot C_{\text{weighted}}(d) + 0.20 \cdot C_{\text{betweenness}}(d) + 0.25 \cdot S(d)
\label{eq:pri}
\end{equation}

\subsection{Component Metrics}

\subsubsection{Degree Centrality $C_{\text{degree}}(d)$}

Captures the interaction burden as the number of drugs interacting with drug $d$, normalized by the maximum possible degree:

\begin{equation}
C_{\text{degree}}(d) = \frac{\text{deg}(d)}{N - 1}
\end{equation}

where $\text{deg}(d)$ is the number of edges connected to drug $d$ and $N$ is the total number of drugs in the network.

\subsubsection{Weighted Degree $C_{\text{weighted}}(d)$}

Represents the severity-weighted sum of all interactions, giving greater importance to contraindicated and major interactions:

\begin{equation}
C_{\text{weighted}}(d) = \frac{\sum_{j \in \Gamma(d)} w_{dj}}{\sum_{i,j} w_{ij}}
\end{equation}

where $\Gamma(d)$ is the set of neighbors of drug $d$ and $w_{dj}$ is the severity weight of the interaction.

\subsubsection{Betweenness Centrality $C_{\text{betweenness}}(d)$}

Quantifies the drug's role in risk propagation pathways across the network, identifying drugs that serve as ``bridges'' between interacting clusters:

\begin{equation}
C_{\text{betweenness}}(d) = \sum_{s \neq d \neq t} \frac{\sigma_{st}(d)}{\sigma_{st}}
\end{equation}

where $\sigma_{st}$ is the number of shortest paths between $s$ and $t$, and $\sigma_{st}(d)$ is the number passing through $d$.

\subsubsection{Severity Profile Score $S(d)$}

Measures the proportion of a drug's interactions that are classified as contraindicated or major severity:

\begin{equation}
S(d) = \frac{|\{j : (d,j) \in E \land \text{severity}(d,j) \in \{\text{Contraindicated}, \text{Major}\}\}|}{|\{j : (d,j) \in E\}|}
\end{equation}

\subsection{Risk Classification Thresholds}

\begin{table}[H]
\centering
\caption{PRI Risk Classification Thresholds}
\label{tab:pri_thresholds}
\begin{tabular}{lll}
\toprule
\textbf{PRI Score} & \textbf{Risk Level} & \textbf{Clinical Action} \\
\midrule
$> 0.5$ & High Risk & Immediate clinical review required \\
$0.3 - 0.5$ & Medium Risk & Close monitoring warranted \\
$< 0.3$ & Lower Risk & Standard monitoring protocols \\
\bottomrule
\end{tabular}
\end{table}

\subsection{Severity Weighting Schema}

The severity weights used in the weighted degree calculation are derived from clinical significance:

\begin{table}[H]
\centering
\caption{Severity Weights for PRI Calculation}
\label{tab:severity_weights}
\begin{tabular}{lr}
\toprule
\textbf{Severity Level} & \textbf{Weight} \\
\midrule
Contraindicated & 4.0 \\
Major & 3.0 \\
Moderate & 2.0 \\
Minor & 1.0 \\
\bottomrule
\end{tabular}
\end{table}

% ============================================================================
\section{Alternative Therapy Suggestion}
% ============================================================================

\subsection{Multi-Objective Drug Recommendation System}

The recommendation system employs a multi-objective optimization framework to identify therapeutic alternatives that balance clinical efficacy with safety. For each candidate alternative drug $d_{\text{alt}}$, we compute a recommendation score:

\begin{equation}
\text{RecScore}(d_{\text{alt}}) = 0.40 \cdot T(d_{\text{alt}}) + 0.35 \cdot \text{SAF}(d_{\text{alt}}) + 0.25 \cdot R(d_{\text{alt}})
\label{eq:recscore}
\end{equation}

\subsection{Score Components}

\subsubsection{Therapeutic Similarity Score $T(d_{\text{alt}})$ (Weight: 0.40)}

Evaluates whether the alternative can serve the same clinical purpose as the original drug. The therapeutic similarity is computed as:

\begin{equation}
T(d_{\text{alt}}) = 0.35 \cdot \text{ATC}(d_{\text{alt}}) + 0.25 \cdot \text{Target}(d_{\text{alt}}) + 0.25 \cdot \text{Disease}(d_{\text{alt}}) + 0.15 \cdot \text{Pathway}(d_{\text{alt}})
\end{equation}

Components include:
\begin{itemize}
    \item \textbf{ATC Code Matching}: Pharmacological class equivalence based on Anatomical Therapeutic Chemical classification
    \item \textbf{Shared Protein Targets}: Similar mechanisms of action through common molecular targets (Jaccard similarity)
    \item \textbf{Disease Indication Overlap}: Therapeutic applicability confirmation
    \item \textbf{Metabolic Pathway Similarity}: Comparable pharmacokinetic profiles
\end{itemize}

\subsubsection{Safety Improvement Score $\text{SAF}(d_{\text{alt}})$ (Weight: 0.35)}

Quantifies the reduction in DDI risk when the alternative is used with other drugs in the current regimen:

\begin{equation}
\text{SAF}(d_{\text{alt}}) = 1 - \frac{\sum_{j \in R} w_{\text{sev}(d_{\text{alt}}, j)}}{\sum_{j \in R} w_{\text{sev}(d_{\text{orig}}, j)}}
\end{equation}

where $R$ is the set of drugs in the current regimen and $w_{\text{sev}}$ is the severity weight.

\subsubsection{Risk Reduction Score $R(d_{\text{alt}})$ (Weight: 0.25)}

Measures the net change in severe interactions and PRI score improvement:

\begin{equation}
R(d_{\text{alt}}) = \frac{\Delta_{\text{contraindicated}} + \Delta_{\text{major}} + \alpha \cdot \Delta_{\text{PRI}}}{\text{max\_possible\_reduction}}
\end{equation}

where $\alpha$ is a scaling factor for PRI improvement contribution.

\subsection{Weight Selection Rationale}

\begin{table}[H]
\centering
\caption{Weight Selection Criteria for Recommendation Score}
\label{tab:weight_rationale}
\begin{tabular}{lcp{8cm}}
\toprule
\textbf{Component} & \textbf{Weight} & \textbf{Rationale} \\
\midrule
Therapeutic Similarity & 0.40 & Highest priority to ensure clinical efficacy is maintained \\
Safety Improvement & 0.35 & Critical for reducing adverse events \\
Risk Reduction & 0.25 & Supporting metric for overall regimen optimization \\
\bottomrule
\end{tabular}
\end{table}

\subsection{ATC Code Matching Levels}

\begin{table}[H]
\centering
\caption{ATC Code Matching Score by Level}
\label{tab:atc_levels}
\begin{tabular}{llc}
\toprule
\textbf{ATC Level} & \textbf{Description} & \textbf{Score} \\
\midrule
5th Level & Same chemical substance & 1.00 \\
4th Level & Same chemical subgroup & 0.80 \\
3rd Level & Same pharmacological subgroup & 0.50 \\
2nd Level & Same therapeutic subgroup & 0.20 \\
1st Level & Same anatomical main group & 0.05 \\
\bottomrule
\end{tabular}
\end{table}

% ============================================================================
\section{Validation Framework}
% ============================================================================

\subsection{Clinical Guidelines Referenced}

The PRI thresholds and recommendation system were validated against established clinical guidelines:

\begin{itemize}
    \item \textbf{AGS Beers Criteria\textsuperscript{\textregistered} 2019}: Potentially inappropriate medications in older adults
    \item \textbf{STOPP/START Criteria v2}: Screening tool for medication review in older persons
    \item \textbf{NHS Scotland Polypharmacy Guidance 2018}: Evidence-based approach to polypharmacy management
\end{itemize}

\subsection{Case Study Design}

Validation employed a clinically realistic 5-drug cardiovascular regimen containing known high-risk interactions, enabling assessment of both risk quantification accuracy and alternative recommendation quality.

% ============================================================================
\section*{References}
% ============================================================================

\bibliographystyle{plain}
\begin{thebibliography}{4}
\bibitem{wishart2018drugbank} Wishart DS, et al. DrugBank 5.0: a major update to the DrugBank database for 2018. \textit{Nucleic Acids Res}. 2018;46(D1):D1074-D1082.
\bibitem{xiong2022ddinter} Xiong G, et al. DDInter: an online drug-drug interaction database. \textit{Nucleic Acids Res}. 2022;50(D1):D1200-D1207.
\bibitem{beers2019} American Geriatrics Society 2019 Beers Criteria Update Expert Panel. American Geriatrics Society 2019 Updated AGS Beers Criteria. \textit{J Am Geriatr Soc}. 2019;67(4):674-694.
\bibitem{stopp2015} O'Mahony D, et al. STOPP/START criteria for potentially inappropriate prescribing in older people: version 2. \textit{Age Ageing}. 2015;44(2):213-218.
\end{thebibliography}

\end{document}
